\paragraph{Results.} The fine grid selected $p=.283$ (clustered misses) and
$.345$ (imperfect clearing). Fresh 512-pair comparisons gave prevention of
72.1\% [69.4, 74.6] and 64.0\% [60.2, 67.5], against 55.9\% and 27.0\% in the
original care-zone configurations, preventing 3.65 and 2.99 acquisitions per run.
Both target checks failed (Table~\ref{tab:gates}): baselines 3.09 [2.85, 3.33]
and 2.85 [2.62, 3.09] (99.167\%) cross the upper and lower edges of the band,
respectively; the clustered descriptive 95\% interval lies inside the band but
cannot replace the wider prespecified check. The clock comparisons change both
the opportunity rule and its calibrated $p$, so their resolved differences do
not isolate a pure clock effect. Under the inherited clock the unfitted
independent $p=1$ boundary had a baseline of 1.70 and prevention of 98.3\%
[97.5, 99.0]; it does not establish a global maximum or an impossibility of
fitting. The independent upper probability ($p=.179$) raised the baseline from
2.78 to 7.29 in the matched 256-seed comparison while prevention point estimates
rose from 66.3\% to 71.0\%; the 4.7-point contrast was unresolved, and this
upward point direction is opposite a simple dilution explanation whose mechanism
is not established.

